% Related Work — NeurIPS framing.
% Three positioning blocks: (a) document KIE landscape with the
% structural-verification gap; (b) verifier-augmented generation in
% NLP; (c) arithmetic-consistency post-processing in invoice
% extraction (the closest published precedent for FOCUS-Σ I3).

\section{Related Work}
\label{sec:related}

\paragraph{Document KIE without verification.}
DONUT~\cite{kim2022donut} ingests an image and autoregressively decodes
structured XML; LayoutLM~\cite{xu2020layoutlm} and its successors
(v2~\cite{xu2021layoutlmv2}, v3~\cite{huang2022layoutlmv3},
LayoutXLM~\cite{xu2021layoutxlm}) align OCR-token text with 2-D
positional embeddings; BROS~\cite{hong2022bros} adds relative spatial
attention; TILT~\cite{powalski2021going}, UDOP~\cite{tang2023unifying},
DocLLM~\cite{wang2023docllm}, and StrucTexT~\cite{li2021structext}
extend with multimodal pre-training objectives.  Pipeline approaches
combine a text-line detector (YOLOv8~\cite{jocher2023yolov8}, DBNet,
or PaddleOCR), a recogniser
(TrOCR~\cite{li2023trocr}, Tesseract), and a learned or
rule-based field assigner.  Across architectures, none of the cited
systems verify their candidate outputs against the receipt's own
arithmetic — they emit one prediction per field per receipt without
post-hoc consistency checks.  Our contribution is a verification layer
that wraps any of these architectures.

\paragraph{Verifier-augmented generation in NLP.}
The pattern of pairing a generator with a downstream verifier has
become standard in NLP for chain-of-thought tasks
(Cobbe et al.~\cite{cobbe2021training} train a verifier to score
solution candidates;  Lightman et al.~\cite{lightman2023verify}
train process- vs outcome-supervised verifiers on math problems;
Madaan et al.~\cite{madaan2023selfrefine} use the generator itself
as iterative critic; Chen et al.~\cite{chen2022codet} verify
generated code against generated tests).
FOCUS-$\Sigma$ is in this lineage but differs in two respects.
\textbf{(i)} The verification target is a single numerical value
bounded by a known arithmetic identity rather than a free-form claim
or program; we exploit this to give exact soundness and completeness
guarantees (Sec.~\ref{sec:focus_sigma}).
\textbf{(ii)} The verifier is itself rule-based (a bounded subset-sum
DP) rather than a learned head, which keeps it
architecture-agnostic and trivially deployable across upstream
extractors.

\paragraph{Arithmetic post-processing in invoice extraction.}
The closest published precedent is in the invoice-extraction
literature, where heuristic checks of \texttt{subtotal + tax = total}
appear as a post-processing step.  Kim et al.~\cite{kim2022donut} cite
``arithmetic-consistency post-processing'' as an ablation;
Hong et al.~\cite{hong2022bros} mention sub-total + tax sanity checks
in their post-process; commercial systems
(Microsoft Form Recogniser~\cite{azure-form-recogniser},
AWS Textract~\cite{aws-textract}) advertise total-field validation
without publishing their rule set.  These approaches are
\emph{keyword-anchored}: they fire only when the \textsf{SUBTOTAL}
and \textsf{TAX} keywords are both present and OCR'd correctly, which
on the canonical SROIE Task-3 split holds for only 38.3\% of
receipts.  FOCUS-$\Sigma$ adds the structural Identity $\mathrm{I}_3$
which fires on the full 87.6\% of receipts that satisfy the
itemisation hypothesis, regardless of which summary keywords the
printer chose to emit.

\paragraph{Subset-sum and combinatorial-verification methods.}
The bounded subset-sum we use is classical knapsack DP
(Karp~\cite{karp1972reducibility}; Bellman~\cite{bellman1957dynamic}).
Its application as a verifier in document AI is, to our knowledge,
new.  Closer methodologically is the program-by-construction
literature
(Solar-Lezama~\cite{solarlezama2008program};
Polozov \& Gulwani~\cite{polozov2015flashmeta}) that synthesises
domain-specific constraints from examples, but those systems
synthesise rules; FOCUS-$\Sigma$ uses a fixed rule (a domain axiom)
and verifies against it.

\paragraph{Receipt benchmarks.}  SROIE~\cite{huang2019icdar} (626
receipts, 4 fields) is the standard single-vendor-type benchmark.
CORD~\cite{park2019cord} (Korean restaurant receipts, 30-field
schema) is the second most-studied receipt corpus.
WildReceipt~\cite{sun2021wildreceipt} extends to in-the-wild
imaging conditions but is rarely used as a fine-tune target.
We evaluate on SROIE Task-3 and CORD-v2; WildReceipt
generalisation is flagged as future work
(Sec.~\ref{sec:limitations}).
